% Supporting detail for readable_full.tex. Keep interpretive caveats in the body.
\section{Data, model, and dictionary details}
\label{app:populations}

\subsection{What is included in the CERT evaluation?}

The positive population is the intersection of the answer key and extracted
session coverage: $70$ user-days from $4$ users on r6.2 and $1{,}309$
from $60$ on r4.2. The corresponding answer-key counts are $99$ rows
from $5$ users and $1{,}883$ from $70$. The four r6.2 users contribute
$46$, $18$, $5$, and $1$ positive days. Claims do not cover incidents
outside the extracted domain.

The original language-model split excludes incident users from adaptation
and assigns the remaining users by a deterministic hash. The validation
comparison contains $405$ benign users on r6.2 and $79$ on r4.2.
Other scored users can have no positive day in the matched domain. These
differences matter: the later r6.2 factorial uses $406$ negative users
($405$ validation users plus one evaluation user with no matched positive),
and r4.2 scoring-time masking includes $89$ never-trained benign users
($79$ validation users plus $10$ others). The classical probe uses its own
random user split (Appendix~\ref{app:classical}).

\subsection{Serialization and adapter training}

The organizational line contains \texttt{week, project, role, b\_unit,
f\_unit, dept, team, itadmin}; the personality line contains
\texttt{O,C,E,A,N}. After \texttt{SESCOUNT}, session lines retain computer
and timing descriptors, work-hour and weekend flags, duration, concurrency,
logon and USB totals, USB duration, file and executable-file totals,
email and attachment statistics, and HTTP categories including job,
cloud, leak, and hacking sites.
The original serializer omitted four hyphenated routing columns:
\texttt{file\_n-to\_usb1}, \texttt{file\_n-from\_usb1},
\texttt{file\_n-file\_act3}, and \texttt{file\_n-disk1}.
The repaired r6.2 reruns restore all four; r4.2 contains one of them.
The original experiments therefore concern the retained field set.
Removing the entire \texttt{DAY} line also removes its week field, so
that condition changes a temporal cue as well as organizational metadata.

All adapters use 4-bit NF4 base weights with double quantization and
bfloat16 computation; LoRA rank $16$, scale $32$, dropout $0.05$;
attention q/k/v/o and feed-forward up/down/gate projections;
sequence limit $2048$; cosine learning-rate decay with $3$ percent
warmup; and zero weight decay. The original Qwen3-8B recipe uses one
epoch, effective batch $8$, and learning rate $1.5\times10^{-4}$.
Qwen2.5-3B uses effective batch $16$ and learning rate $2\times10^{-4}$.
The 8B factorial uses batch $16$ across four GPUs while retaining its
$1.5\times10^{-4}$ rate. Recipes are held fixed across input conditions
within each model; they are not identical across model families.

The repaired 3B runs use one epoch on all $287{,}827$ r4.2 training days
or a $300{,}000$-day cap from approximately $1.25$ million r6.2 days,
with adapter seeds $42$ and $43$. TWOS uses three epochs on its smaller
benign corpus. The factorial and LANL runs use a $300{,}000$-example
training cap and a $2{,}000$-example validation cap, seed $42$.

Scores average next-token negative log-likelihood over attention-valid
targets; padding is excluded. When full fp32 logits would exceed
$28$ GiB, scoring computes cross-entropy in bounded slices from the
hidden states and output embedding. This changes memory use, not the
score's mathematical definition. Input removal also changes which tokens
enter that average, one reason the factorial is a pipeline-level effect.

\subsection{Dictionary fitting and feature selection}

Token deltas come from the residual stream after the audited layer,
in evaluation-example order. For coordinatewise mean $\mu$ and standard
deviation $\sigma$, the encoder and decoder are
\[
 z_t=\operatorname{TopK}_k\!\left(\operatorname{ReLU}
 (W_e((\delta_t-\mu)/\sigma)+b_e)\right),\qquad
 \widehat\delta_t=\sigma\odot Dz_t+\mu.
\]
The decoder is untied and has no bias in standardized coordinates.
Training uses Adam, learning rate $10^{-3}$, $20$ epochs, batches of
$1024$--$4096$, and at most two million token rows per layer. The
original fit retains all positive-example rows and subsamples benign
rows; benign-only dictionaries are separately identified.

Here $m$ is the number of dictionary columns. The core search varies
$m/d\in\{2,4\}$, $k\in\{4,8\}$, and layers $18,26,34$ on r6.2
or $18,26$ on r4.2. The headline r6.2 choice is layer $18$, $m/d=4$,
$k=8$: standardized reconstruction MSE $0.054$, mean active count
$7.98$, and $1.0$ percent dead features. The native r4.2 choice is
layer $26$, $m/d=2$, $k=4$: MSE $0.093$, active count $4.0$,
and $1.4$ percent dead features. These configurations were selected
using exploratory intervention results on the full positive population.

The candidate set contains the five largest positive-minus-benign mean
activation gaps. The active control contains five other features with
smallest absolute gap among those active on at least $0.002$ of token
rows. Session-restricted selection recomputes gaps on \texttt{SES} rows
and excludes features with more than $0.5$ profile mass on positive rows.
Held-out analyses reselect feature identities, but inherit the dictionary,
layer, size, sparsity, context menu, and edit grid unless explicitly refitted.

\section{Detector measurements and input controls}
\label{app:detectors}

\subsection{Original baseline comparison}

Each original fold (seed $42$) holds out one malicious user and samples
$800$ benign users. The language model is fixed across folds and has
trained on many of those benign users; baseline detectors retrain with
test users excluded. Day metrics cover all test rows; user metrics use
the maximum day score. Table~\ref{tab:readable_detectors} reports the
fold averages. Deep SVDD is deep support vector data description
\citep{ruff2018deepsvdd}; GRU and LSTM AE are gated recurrent unit and
long short-term memory autoencoders; IF is Isolation Forest
\citep{liu2008isolation}. Their distinct training exposure prevents
interpreting the LM's higher ROC as a fair generalization advantage.

\begin{table}[htbp]
\centering\small
\begin{tabular}{llrrrr}
\toprule
Dataset & Detector & Day ROC & Day AP & User AP & Mean rank \\
\midrule
r6.2 & Adapted LM & $.953$ & $.000755$ & $.054$ & $24.0$ \\
 & Deep SVDD & $.628$ & $.01155$ & $.211$ & $83.2$ \\
 & GRU AE & $.766$ & $.00572$ & $.081$ & $24.8$ \\
 & LSTM AE & $.768$ & $.00207$ & $.057$ & $24.2$ \\
 & IF & $.713$ & $.000211$ & $.013$ & $153.0$ \\
\midrule
r4.2 & Adapted LM & $.964$ & $.01345$ & $.101$ & $26.4$ \\
 & Deep SVDD & $.743$ & $.03372$ & $.382$ & $53.4$ \\
 & GRU AE & $.696$ & $.02544$ & $.124$ & $86.6$ \\
 & LSTM AE & $.714$ & $.02364$ & $.120$ & $92.1$ \\
 & IF & $.715$ & $.000254$ & $.008$ & $186.4$ \\
\bottomrule
\end{tabular}
\caption{Original fold-aligned comparison. Rank is the malicious user's
position among $801$ users (smaller is better). AP is average precision,
called PR-AUC in the result artifacts. Baseline user ROC is unavailable
in this harness; adapted-LM fold averages are approximately $.971$ and
$.968$ from the tracked rows. These are not the pooled values below.}
\label{tab:readable_detectors}
\end{table}

Base-model NLL has day ROC $0.461$ on r6.2 and $0.581$ on r4.2 in
the original comparison. Thus adaptation changes the ranking substantially,
without establishing whether the acquired information is appropriate.
In the full scored pool, restricting benigns to validation users changes
day ROC from $0.954$ to $0.545$ on r6.2 and $0.959$ to $0.668$ on
r4.2; user ROC changes from $0.970$ to $0.546$ and $0.969$ to $0.565$.
The restricted day/user APs are $0.0005/0.013$ and $0.1023/0.521$.
These APs have different class prevalence from the fold comparison and
should not be treated as a direct cross-protocol performance gain.

\subsection{Within-user ranking, masking, and swaps}

Mean within-user ROC is $0.725$ on r6.2, with user-bootstrap interval
$[0.654,0.795]$ and all four users above chance; the dominant user has
$0.768$. On r4.2 the mean is $0.669$, interval $[0.619,0.717]$,
with $52$ of $60$ users above chance. These compare each user's
positive and own-benign days, not seen versus unseen identities.

Scoring-time masking uses a fixed adapter and a fixed unseen comparison.
The r4.2 pool contains $42{,}468$ days from $149$ users, including the
$89$ benign users described above. Removing personality changes day
ROC from $0.670$ to $0.797$ (paired change $+0.127$, interval
$[0.061,0.193]$), day AP from $0.069$ to $0.112$, and user ROC
from $0.666$ to $0.775$ (change $+0.108$, interval $[0.038,0.179]$).
Further masking organizational identifiers gives day change $+0.063$
with interval $[-0.003,0.130]$. On r6.2, original/no-personality/
no-profile day ROC is $0.547/0.575/0.471$; all paired change intervals
contain zero. Masking intervals use $10{,}000$ paired bootstrap draws,
resampling users within both label strata. Scoring-time masking and
the later removal-and-retraining experiment are distinct interventions.

The swap probe replaces either the positive day's profile or its session
content with an unseen validation user's values, keeping the complementary
part fixed. Mean signed NLL changes are $+0.155$ versus $+0.072$ on
r6.2 and $+0.180$ versus $+0.102$ on r4.2. Both profile swaps
increase mean NLL more. These replacements can disrupt natural
profile--behavior combinations; they measure sensitivity, not an
isolated familiarity mechanism or a causal feature's semantic identity.

\section{Exact interventions and uncertainty}
\label{app:intervention}

\subsection{Repair and ablation}

Let $F$ be the candidate or control feature set and
$T_F=\{t:\sum_{f\in F}z_{t,f}>0\}$ the active receiver positions.
For a donor $d$, the prototype $p_f^d$ is the mean donor activation
over positions active for $F$, falling back to all donor positions if
that set is empty. For $f\in F$ and $t\in T_F$, repair sets
\[
 z'_{t,f}=(1-\alpha)z_{t,f}+\alpha p_f^d,
 \qquad \alpha\in\{0.25,0.5,0.75,1\}.
\]
Other coefficients are unchanged. At every token, the patched delta is
$\delta'_t=\sigma\odot Dz'_t+\mu$; a forward hook adds
$\delta'_t-\delta_t$ to the adapted residual state. Thus both arms
substitute a sparse reconstruction even away from edited positions.
That shared baseline need not cancel after nonlinear downstream layers.

For receiver $r$ and donor class $b\in\{\mathrm{ben},\mathrm{anom}\}$,
the recorded outcome is
\[
 \Delta^F_b(r)=\min_{d\in\mathcal D_b(r),\,\alpha}
 \{s(x_r;F,d,\alpha)-s(x_r)\}.
\]
Negative values mean score reduction. For the same complete receivers
$\mathcal R$, repair is
\[
 \tau=\frac1{|\mathcal R|}\sum_{r\in\mathcal R}
 \left[(\Delta^S_{\mathrm{anom}}(r)-\Delta^S_{\mathrm{ben}}(r))
 -(\Delta^C_{\mathrm{anom}}(r)-\Delta^C_{\mathrm{ben}}(r))\right].
\]
This is maximal achievable repair under the permitted search, not a
donor-averaged causal effect. Donor pools are capped at $16$ per receiver,
drawn with a fixed random seed; malicious donors are scarcer (median
$1$ on r6.2, mean $13$ on r4.2), creating unequal search opportunities.

Ablation replaces the edited coordinates by $(1-\alpha)z_{t,f}$.
For positive or matched-benign receiver $q$, let
$e_F(q)=\min_\alpha\{s(x_q;F,\alpha)-s(x_q)\}$. Over complete
positive/benign pairs $(r,b)$, the ablation contrast is
\[
 \nu=\overline{[e_S(b)-e_S(r)]-[e_C(b)-e_C(r)]}.
\]
Positive $\nu$ means more positive-specific score reduction from
ablating the selected set than the control. It does not prove strict
necessity of the features outside this operation.

Donor and receiver pairs match role, department$\times$role,
project$\times$role, team, or department as specified by the context.
All donors and benign receiver matches exclude the positive receiver's
user. Benign matches come from never-trained evaluation users.
r4.2 repair has $1034$--$1157$ complete receivers across its contexts.
r6.2 team has no complete malicious-donor contrast and is unavailable.
Every contrast uses complete paired outcomes, not separate arm means
with different missingness.

\subsection{What the control checks establish}

The initial active controls are not matched in edit size. Typical best-row
repair norms are approximately $71$ versus $20$ (selected versus control)
on r6.2 and $112$--$124$ versus $11$--$13$ on r4.2. Within a feature
set, benign- and malicious-donor norms differ by only $0.2$ percent on
r6.2 and at most $6$ percent on r4.2. These observations help diagnose
the contrast but do not eliminate magnitude-related confounding.
The reconstruction component has mean per-token norm approximately
$7$ and $37$, respectively.

Controls additionally matched on activation frequency, magnitude, and
decoder norm reduce the selected/control norm ratio to $1.33$ on r6.2
and $1.79$ on r4.2. The rerun results are:
\begin{center}\small
\begin{tabular}{lcc}
\toprule
 & Repair $\tau$ [95\% interval] & Ablation $\nu$ [95\% interval] \\
\midrule
r6.2 & $.0059$ [$-.0018,.0096$] & $.0609$ [$ .0074,.0843$] \\
r4.2 & $.00168$ [$ .00105,.00228$] & $.0012$ [$-.0016,.0038$] \\
\bottomrule
\end{tabular}\end{center}
The intervals are user-cluster bootstrap intervals; r6.2 has only four
clusters. The r4.2 repair and r6.2 ablation contrasts retain intervals
above zero. This is partial calibration, not exact equal-dose testing.

Replacing best-donor search with the median of each donor's best-strength
outcome yields positive headline repair contrasts of $0.0005$ (r6.2)
and $0.0006$ (r4.2). Dividing each best-row score change by its edit
norm gives r4.2 positive point estimates in three of four contexts
(team $+0.0008$ per norm unit), but all context intervals include zero.
r6.2 repair does not survive this rescaling (three usable user clusters).
The rescaling assumes approximately linear response and is a sensitivity
check, not an alternative identification proof.

Ablation sizes also differ: selected/control norms approximately
$45/21$ on r6.2 and $165$--$169$/$13$--$15$ on r4.2.
Selected features have fewer active tokens on benign matches
($0.2$--$0.8$ versus $3.0$ on r6.2; about $2$ versus $4.4$ on
r4.2). The receiver-class contrast is therefore partly a comparison
between small and large edits. The matched-control and replication
checks are material qualifications to its interpretation.

\subsection{Uncertainty and selection}

Primary intervention intervals use $10{,}000$ draws of malicious users
with replacement and recompute the pooled complete-case contrast.
Users with more eligible days retain greater weight in the point
estimate; clustering does not make it a mean of equally weighted users.
Supplementary procedures use $4000$ receiver/pair draws or bootstrap a
mean of user means. The figures use the primary user-cluster result.
All contexts are reported without a multiple-comparison correction;
a strongest-context example is not confirmation of every context.
Intervals condition on trained models and selected configurations,
and four-user intervals are descriptive.

\section{Transfer and replication details}
\label{app:robustness}

\subsection{Literal transfer and geometric alignment}
\label{app:transfer}

Transferring the r6.2 dictionary and feature sets directly to all $1309$
r4.2 positive receivers gives negative repair contrasts in every audited
configuration. The best values in the three transferred sparsity settings
are $-0.00094$, $-0.00105$, and $-0.00113$.
Native r4.2 rediscovery at layer $26$ instead gives positive repair
in all four contexts: team $0.00142$, role $0.00111$,
department$\times$role $0.00107$, department $0.00098$, with primary
intervals above zero. Native ablation is positive with intervals above
zero for department$\times$role ($0.0029$) and role ($0.0021$),
while team and department intervals cross zero (Figure~\ref{fig:effects}).
A local session autoencoder evaluated at the same user-day unit gives
best repair contrasts $0.0011$ on r6.2 and $0.00091$ on r4.2.
Its reconstruction-error score has different units from LM NLL, so
these numbers do not quantify a relative causal advantage.

Feature alignment uses absolute cosine after undoing standardization:
compare normalized $\sigma\odot D_{:,f}$. Each selected source direction
is matched to the best target direction. The reference is the median
best-match similarity across the whole source dictionary, not a formal
statistical null test.
\begin{center}\small
\begin{tabular}{lcc}
\toprule
Comparison & Selected best matches & Dictionary reference \\
\midrule
r6.2, three seed pairs & $.807$--$.856$ & $.720$--$.750$ \\
r4.2, three seed pairs & $.837$--$.946$ & $.814$--$.818$ \\
Across datasets, layers 18/26 & $.218$--$.236$ & $.224$--$.405$ \\
Across datasets, layer 18 & $.295$--$.652$ & $.370$--$.963$ \\
\bottomrule
\end{tabular}\end{center}
r4.2's layer-18 comparison is away from its headline layer-26 mechanism.
Fresh r6.2 dictionary seeds also yield positive repair in every finite
context (role $0.0071$ and $0.0019$, versus original $0.0068$).
This tests dictionaries on the same cached representations.

\subsection{Held-out users and benign-only dictionaries}

r4.2 discovery/confirmation splits $60$ malicious users into $30/30$
($671/638$ positive days). Four of five selected features agree with
full-population selection. Held-out repair is $0.00034$--$0.00076$;
only department has a primary interval above zero ($0.00076$,
$[0.00035,0.00121]$, $21/29$ contributing users positive).
Held-out ablation is directionally positive for role and
department$\times$role; team does not replicate a positive effect.

With the original r6.2 selection, project$\times$role ablation is
positive for all four users ($0.019$--$0.103$); pooled repair is largely
driven by the $46$-day user. In leave-one-user-out selection, the three folds retaining the
dominant user select nearly identical sets; excluding that user changes
the set completely. The new set repairs the dominant user's held-out
days ($0.0037$--$0.0038$), while effects for low-activity users are
near zero. Held-out ablation is positive on the two other multi-day
users (up to $0.043$) but negative on the dominant-user fold.
Intervals within a single held-out user describe variability over
that user's days, not independent-user replication.

With dictionaries fitted only to benign deltas, r4.2 attribution remains
at least $94$ percent session mass, with four features at $100$ percent.
Held-out ablation is $+0.004$--$+0.013$ across contexts
($528$--$638$ complete pairs), but repair is $-0.0002$--$-0.0004$.
r6.2 remains at least $99$ percent profile mass in every fold;
repair is positive in one held-out fold, while ablation is positive
on the dominant-user fold. Dictionary choice changes which intervention
and user combination replicates.

\subsection{Population, new adapters, and secondary features}

For the fixed r4.2 model and dictionary, positive-user subsamples of
$4,10,20,30,60$ (ten draws each) give mean top-five profile share
about $0.25$ at four users and $0.16$--$0.20$ otherwise. This is a
selection-stage test conditional on a full-population dictionary.

The repaired 3B adapters use fresh deltas at layers $12,18,24$ and
benign-only dictionaries. On r4.2, all five features at the best layer
have at least $94$ percent session mass under both adapter seeds;
unseen-user day ROC is $0.678/0.718$. Ablation is $+0.00218/0.00215$
on $1301$ complete pairs per seed, with repair about $-0.0004$ each.
On r6.2, three or four of five layer-24 features are almost entirely
profile-bound; shallower configurations can be behavioral. Day ROC is
$0.414/0.738$; repair is $+0.0175$ in one seed and near zero in the other.
These runs jointly change model, serializer, and training realization.

r6.2 session-restricted leave-one-user-out selection recovers two recurring
features across three of four folds, with near-zero profile mass.
Held-out repair on the dominant user is $+0.0027$--$+0.0029$, around
three-quarters of the unrestricted profile-family effect. Ablation
replicates partially, including the dominant-user role context where
the unrestricted profile-family result was negative. This shows a
secondary session-concentrated family within the same model.

\section{The two real-data studies}
\label{app:realdata}

TWOS contains $24$ real users over five days with real IPIP-50 personality
questionnaires. The serialization yields $18{,}048$ ten-minute windows,
$138$ positive ($101$ masquerade and $37$ traitor). The benign training
pool is approximately $7000$ windows from $12$ users never masqueraded;
future traitors contribute only pre-dismissal benign days. Both adapter
seeds and audited layers $18/24$ yield at least $0.983$ session mass
for each selected feature. Within-user ROC averages $0.783$ over the
$16$ attacked users with logged windows. The quoted interventions
($\nu=+0.0019$, $\tau=-0.0018$) use seed $42$, layer $24$,
team matching, same-user exclusion, and $138$ receivers. Dictionaries
are benign-only but feature selection uses the full positive pool;
the quoted point estimates are exploratory and are not accompanied
here by independent-user confirmation.

LANL contains approximately $1.05$ billion authentication events over
$58$ days. Keeping all $104$ red-team users and a one-in-ten sample
of other users yields approximately $1.62$ million windows, $432$
positive. A window contains up to $32$ events within a user-day and
is positive if it includes a red-team event. Five deterministic MD5
user folds assign fold $4$ to unseen evaluation; benign windows from
folds $0$--$3$ supply training and validation. Evaluation keeps all
positives and hash-samples benign windows, disjoint from the training
and validation windows. Actual scored counts, rather than target sample
budgets, are $22{,}563$ seen-fold windows ($234$ positive) and
$30{,}410$ unseen-fold windows ($198$ positive). A seen-fold user can
have benign training windows even if another window is red-team positive.

Token-length checks cover all $52{,}973$ scored windows under the full
and host-anonymized conditions: full maximum $1503$, 99th percentile
$1408$, positive maximum $1417$; host-anonymized maximum $1311$,
99th percentile $1263$, positive maximum $1212$. All are below $2048$.
This check should not be expanded into a claim about every raw or
generated window in all conditions. The host-constancy audit covers
$9{,}414{,}753$ training events and $1180$ users with at least $20$
events: modal source-host share has median $0.54$, quartiles
$0.40/0.61$, mean $0.51$; modal destination-host share has median
$0.16$. The full versus host-anonymized unseen AP is approximately
$0.0133$ versus $0.0129$, at prevalence $0.00651$. These values are
above prevalence but show no AP improvement from host masking.

\section{Classical and train-matched comparisons}
\label{app:classical}

\subsection{The separate classical probe}

The classical probe aggregates numeric session columns to user-days:
profile columns take the first value; behavioral columns take means,
with session count added. It includes $12$ profile and $230$ behavioral
columns on r6.2, $11$ and $110$ on r4.2. It randomly holds out
$399/94$ benign users, respectively (seed $42$), and scores their days
plus all matched positive days. This differs from the LM's hash split.
Training-day counts are $1{,}254{,}025/284{,}322$ and evaluation-day
counts $138{,}138/33{,}789$.

Features are standardized on training data. Isolation Forest uses
$300$ trees; the one-class SVM uses $\nu=0.05$, scaled kernel width,
and a $6000$-row training subsample; principal-component analysis (PCA)
uses $10$ components with reconstruction norm as score. The input-removal
variants refit each model; SVM subsamples need not be identical between
variants. Permutation importance is one random permutation per column,
with negative ROC drops clipped to zero before computing group shares.
These shares are descriptive and differ from sparse activation mass.

\begin{center}\small
\begin{tabular}{llrrrr}
\toprule
 & Model & Profile share & Profile in top 5 & Full ROC & Removal change \\
\midrule
r6.2 & IF & $.040$ & $0$ & $.6408$ & $+.0030$ \\
 & One-class SVM & $.065$ & $0$ & $.4914$ & $+.0225$ \\
 & PCA & $.041$ & $0$ & $.6164$ & $+.0220$ \\
r4.2 & IF & $.017$ & $0$ & $.8294$ & $-.0084$ \\
 & One-class SVM & $.124$ & $1$ & $.6595$ & $-.0223$ \\
 & PCA & $.024$ & $0$ & $.8382$ & $+.0089$ \\
\bottomrule
\end{tabular}\end{center}
Profile reliance is small in the r6.2 probe. r4.2 is mixed: personality
trait O enters the SVM's top five, and removal lowers two detectors'
ROC. Thus the r6.2 pattern cannot be asserted unchanged on r4.2.

\subsection{Factorial intervals and individual users}
\label{app:factorial}

Each arm uses the same capped benign corpus and input condition during
adaptation and scoring. The audit pool includes all evaluation-user days
and a seeded $12$ percent sample of validation days. Its fold re-score
ranks each malicious user against the same $406$ negative users.
User AUC with one positive is a measure of that user's position relative
to negatives, not the probability of detecting that user at a chosen
alert threshold. Table~\ref{tab:readable_factorial} retains the individual
results behind the body averages.

\begin{table}[htbp]
\centering\small
\begin{tabular}{llcccc}
\toprule
 & Condition & ACM2278 & CDE1846 & CMP2946 & MBG3183 \\
\midrule
3B & Full & $.995$ (3) & $.803$ (81) & $.966$ (15) & $.941$ (25) \\
 & No personality & $.916$ (35) & $.926$ (31) & $.924$ (32) & $.845$ (64) \\
 & No profile & $.993$ (4) & $.993$ (4) & $.993$ (4) & $.845$ (64) \\
 & Shuffle & $.894$ (44) & $.488$ (209) & $.520$ (196) & $.722$ (114) \\
\midrule
8B & Full & $.919$ (34) & $.155$ (344) & $.527$ (193) & $.527$ (193) \\
 & No personality & $.931$ (29) & $.872$ (53) & $.911$ (37) & $.911$ (37) \\
 & No profile & $.867$ (55) & $.978$ (10) & $.951$ (21) & $.818$ (75) \\
 & Shuffle & $.345$ (267) & $.241$ (309) & $.574$ (174) & $.778$ (91) \\
\bottomrule
\end{tabular}
\caption{User ROC and rank in parentheses, out of $407$ users per fold.
All four users rank in the top fifth after profile removal in both
models. This is a ranking result, not a demonstrated alert recall.}
\label{tab:readable_factorial}
\end{table}

The $95$ percent bootstrap intervals for user ROC are:
\begin{center}\small
\begin{tabular}{lcc}
\toprule
Condition & 3B & 8B \\
\midrule
Full & $[.844,.982]$ & $[.248,.821]$ \\
No personality & $[.865,.925]$ & $[.882,.926]$ \\
No profile & $[.882,.993]$ & $[.842,.964]$ \\
Shuffle & $[.504,.808]$ & $[.293,.676]$ \\
\bottomrule
\end{tabular}\end{center}
Paired profile-removal changes are $+.030$ $[-.065,.142]$ and
$+.371$ $[.067,.690]$. Resampling uses the same four malicious-user
indices across contrasted arms and fixes the negative cohort. Bootstrap
tail fractions in the raw logs are not additional training replications
or a four-user significance guarantee.

In the separately fitted full/shuffle dictionaries, selected profile
fractions are:
\begin{center}\small
\begin{tabular}{lccccc}
\toprule
Model, layer, condition & Feature 1 & 2 & 3 & 4 & 5 \\
\midrule
3B, 12, full & $1$ & $.853$ & $.997$ & $.026$ & $.992$ \\
3B, 12, shuffle & $1$ & $.920$ & $1$ & $1$ & $.985$ \\
8B, 24, full & $.053$ & $.041$ & $.095$ & $1$ & $1$ \\
8B, 24, shuffle & $.085$ & $.095$ & $.257$ & $.998$ & $1$ \\
\bottomrule
\end{tabular}\end{center}
These are the capped factorial adapters, not the original layer-18
8B audit; each has its own selected features.

The exploratory novelty check compares each positive user with $3590$
benign training users. In the table's user order, counts of training
profiles within Big-Five $L_1$ distance $10$ are $3,12,11,0$;
minimum distances are $9,7,6,12$; exact role/department/team matches
are $29,13,22,4$. Corresponding full-8B day ROC values are
$0.340,0.191,0.412,0.653$. The nearest-distance and neighborhood-count
rankings differ, and neither perfectly orders user ROC. This four-case
post hoc comparison cannot establish a general novelty mechanism.
